% Direct-authority diplomatic French for printed p.183.
\LayerTarget{tohoku:proposition:3.10.2}
\noindent\textsc{Proposition 3.10.2.}\SourcePage{183} \textit{Supposons l'espace localement
compact $X$ dénombrable à l'infini. Alors les homomorphismes (3.10.2) (où la
limite projective est prise suivant l'ensemble filtrant croissant des parties
ouvertes relativement compactes $U$ de $X$) sont surjectifs. Si $p>1$, pour
qu'il soit bijectif, il suffit que pour tout $U$ ouvert relativement compact,
en existe un autre $V\supset U$ tel que pour tout ouvert relativement compact
$W$ contenant $V$, l'image dans $H^{p-1}(U,F)$ de $H^{p-1}(W,F)$ par
l'homomorphisme de restriction soit identique à celle de $H^{p-1}(V,F)$. Si
$p=1$, pour que l'homomorphisme
$H^1(X,F)\longrightarrow\varprojlim_U H^1(U,F)$ soit bijectif, il suffit qu'on
puisse munir les $H^0(U,F)$ de topologies de groupes topologiques métrisables
complets telles que les homomorphismes de restriction soient continus, et que
pour tout ouvert relativement compact $U$, en existe un autre $V\supset U$
tel que pour tout ouvert relativement compact $W$ contenant $V$, l'image de
$H^0(W,F)$ dans $H^0(U,F)$ soit dense dans celle de $H^0(V,F)$.}

\LayerTarget{tohoku:remark:proposition-3.10.2:approximation}
(Bien entendu, on pouvait dans cet énoncé remplacer les ouverts relativement
compact par des compacts). Cette proposition, qui se démontre par un procédé
d'approximation à la Mittag-Leffler, est par exemple essentielle dans la
démonstration des théorèmes fondamentaux sur les variétés de Stein [5]. Pour
$p=1$, elle reste vraie si $F$ est un faisceau de groupes non nécessairement
abéliens, sous la forme : si la condition d'approximation énoncée (dans le cas
$p=1$) est vraie, tout élément du premier membre de (3.10.2) dont l'image est
l'élément neutre, est l'élément neutre.

\TohokuPartBoundary{185}
